% Part: incompleteness
% Chapter: incompleteness-provability
% Section: fixed-point-lemma

\documentclass[../../../include/open-logic-section]{subfiles}

\begin{document}

\olfileid{inc}{inp}{fix}

\olsection{स्थिरबिंदू उपपत्ती}

\begin{explain}
स्थिरबिंदू उपपत्ती सांगते की कोणत्याही
!!{formula}~$!B(x)$ साठी असे !!a{sentence}~$!A$ असते की
$\Th{T} \Proves !A \liff !B(\gn{!A})$, मात्र
$\Th{T}$ हा~$\Th{Q}$ चा विस्तार असला पाहिजे.
खोटारड्याच्या वाक्यासाठी $!A$ हे~“$\gn{!A}$ असत्य आहे”
या विधानाशी (उपपत्ती~$\Th{T}$ मध्ये सिद्ध होईल अशा
रीतीने) सममूल्य हवे; म्हणजेच $\Gn{!A}$ ही एका असत्य
!!{sentence} ची गोडेल संख्या आहे असे त्यात म्हटलेले हवे.
सिद्धतेची कल्पना समजून घेण्यासाठी $!A$ विषयी क्वाइनच्या
अनौपचारिक मांडणीशी तुलना करू: “उद्धृत रूप आधी जोडल्यास
असत्य विधान मिळते” उद्धृत रूप आधी जोडल्यास असत्य विधान
मिळते. एखादा शब्दसमूह घेऊन त्याच्या आधी त्याचेच उद्धृत
रूप जोडून वाक्य तयार करण्याच्या क्रियेला त्या शब्दसमूहाचे
\emph{विकर्णीकरण} म्हणू शकतो; तयार झालेला परिणाम हे
त्याचे विकर्णरूप. म्हणून “उद्धृत रूप आधी जोडल्यास असत्य
विधान मिळते” या शब्दसमूहाचे विकर्णरूप म्हणजे “उद्धृत रूप
आधी जोडल्यास असत्य विधान मिळते” उद्धृत रूप आधी
जोडल्यास असत्य विधान मिळते. आता लक्षात घ्या की क्वाइनचे
खोटारड्याचे वाक्य हे “असत्य विधान मिळते” याचे नव्हे,
तर “उद्धृत रूप आधी जोडल्यास असत्य विधान मिळते” याचे
विकर्णरूप आहे. म्हणजेच खोटारड्याचे वाक्य मिळवण्यासाठी
ज्याचे विकर्णीकरण करतो, त्या गुणधर्मातच विकर्णीकरण
समाविष्ट असते!{}

अंकगणिताच्या भाषेत, एक मुक्त चर असलेल्या
!!a{formula} ची गोडेल संख्या काढून, त्या चराच्या जागी
त्या संख्येचा प्रमाण संख्यांक बसवून त्याचे उद्धृत रूप
बनवतो. $!E(x)$ चे विकर्णरूप हे $!E(\num{n})$ आहे,
जिथे $n = \Gn{!E(x)}$. (यापुढे
$\num{\Gn{!E(x)}}$ याचे संक्षिप्त रूप
$\gn{!E(x)}$ असे लिहू.) म्हणून $!B(x)$ हा
“असत्य आहे” हा गुणधर्म असेल, तर “स्वतःचे उद्धृत रूप
आधी जोडल्यास असत्य विधान मिळते” याचा अर्थ
“स्वतःच्या विकर्णरूपाच्या गोडेल संख्येला लागू केल्यावर
असत्य विधान मिळते” असा होईल. !!{formula} ची गोडेल
संख्या~$n$ दिल्यावर त्याच्या विकर्णरूपाची गोडेल संख्या
काढणाऱ्या $\fn{diag}(n)$ या फलनासाठी
$\Obj{diag}$ असे चिन्ह आपल्याकडे असेल, तर
$!E(x)$ हे $!B(\Obj{diag}(x))$ असे लिहिता येईल.
क्वाइनच्या खोटारड्याच्या वाक्याचे रूप मग याचे
विकर्णीकरण, म्हणजे $!E(\gn{!E(x)})$ किंवा
$!B(\Obj{diag}(\gn{!B(\Obj{diag}(x))}))$, असे असेल.
अर्थात, आता $!B(x)$ हा इतर कोणताही गुणधर्म असला,
तरी हीच रचना चालेल. अपूर्णता प्रमेयासाठी $!B(x)$
हा “$x$ हे~$\Th{T}$ मध्ये !!{derivable} नाही” हा
गुणधर्म घेऊ. तेव्हा $!E(x)$ चे वर्णन असे होईल:
“स्वतःच्या विकर्णरूपाच्या गोडेल संख्येला लागू केल्यावर
$\Th{T}$ मध्ये !!{derivable} नसलेले
!!a{sentence} मिळते.”

$\Th{T}$ मध्ये ही रचना आकारिक रीतीने मांडण्यासाठी
$\fn{diag}$ चे आकारिकीकरण करावे लागेल.
$\fn{diag}(n)$ हे संगणनक्षम, किंबहुना आदिम पुनरावर्ती
फलन आहे: $n$ ही~$!E(x)$ या सूत्राची गोडेल
संख्या असेल, तर $\fn{diag}(n)$ हे
$!E(\gn{!E(x)})$ ची गोडेल संख्या देते.
(आठवा, $\gn{!E(x)}$ हा~$!E(x)$ च्या गोडेल
संख्येचा प्रमाण संख्यांक, म्हणजे
$\num{\Gn{!E(x)}}$, आहे.) $\Obj{diag}$ हे
$\fn{diag}$ चे निरूपण करणारे $\Th{T}$ मधील
फलनचिन्ह असेल, तर $!A$ म्हणून
$!B(\Obj{diag}(\gn{!B(\Obj{diag}(x))}))$
हे सूत्र घेता येईल. लक्षात घ्या:
\begin{align*}
\fn{diag}(\Gn{!B(\Obj{diag}(x))}) & =
\Gn{!B(\Obj{diag}(\gn{!B(\Obj{diag}(x))}))} \\
& = \Gn{!A}.
\end{align*}
$\Th{T}$ पुढील विधान !!{derive} करू शकते असे धरल्यास,
\[
\Obj{diag}(\gn{!B(\Obj{diag}(x))}) = \gn{!A},
\]
त्याला $!B(\Obj{diag}(\gn{!B(\Obj{diag}(x))}))
\liff !B(\gn{!A})$ हेही !!{derive} करता येते.
पण डावी बाजू व्याख्येनुसार~$!A$ आहे.

अर्थात, सामान्यतः $\Obj{diag}$ हे
$\Th{T}$ मधील फलनचिन्ह नसेल; आणि
$\Th{Q}$ मधील तर नक्कीच नाही. पण
$\fn{diag}$ संगणनक्षम असल्यामुळे ते
$\Th{Q}$ मध्ये एखाद्या
$!D_{\fn{diag}}(x,y)$ सूत्राने
\emph{निरूपित} करता येते. म्हणून
$!B(\Obj{diag}(x))$ असे लिहिण्याऐवजी
$\lexists[y][(!D_{\fn{diag}}(x,y) \land !B(y))]$
असे लिहू शकतो. याखेरीज वर रेखाटलेली सिद्धता तशीच
चालते; किंबहुना ती आधीच~$\Th{Q}$ मध्ये चालते.
\end{explain}

\begin{lem}
\ollabel{lem:fixed-point} $!B(x)$ हे~$x$ हे एक मुक्त चर
असलेले कोणतेही सूत्र असू द्या. मग असे वाक्य~$!A$
असते की $\Th{Q} \Proves
!A \liff !B(\gn{!A})$.
\end{lem}


\begin{proof}
$!B(x)$ दिल्यावर $!E(x)$ हे
$\lexists[y][(!D_{\fn{diag}}(x,y) \land !B(y))]$
सूत्र मानू, आणि $!A$ हे त्याचे विकर्णरूप,
म्हणजे $!E(\gn{!E(x)})$, मानू.

$!D_{\fn{diag}}$ हे $\fn{diag}$ चे निरूपण
करते आणि $\fn{diag}(\Gn{!E(x)}) = \Gn{!A}$
असल्यामुळे $\Th{Q}$ पुढील विधाने
!!{derive} करू शकते:
\begin{align}
  & !D_{\fn{diag}}(\gn{!E(x)}, \gn{!A}) \ollabel{repdiag1} \\
  & \lforall[y][(!D_{\fn{diag}}(\gn{!E(x)},y) \lif
  \eq[y][\gn{!A}])]. \ollabel{repdiag2}
\end{align}
आता $\Th{Q} \Proves !A \liff !B(\gn{!A})$
हे दाखवू. केवळ तर्कशास्त्र आणि $\Th{Q}$
मध्ये !!{derivable} असलेली तथ्ये वापरून
अनौपचारिक युक्तिवाद करू.

प्रथम~$!A$, म्हणजे $!E(\gn{!E(x)})$, खरे
मानू. $!E(x)$ च्या व्याख्येकडे परत पाहिल्यास,
$!E(\gn{!E(x)})$ हेच पुढील विधान आहे:
\[
\lexists[y][(!D_{\fn{diag}}(\gn{!E(x)},y) \land !B(y))].
\]
अशा~$y$ चा विचार करा. आता
$!D_{\fn{diag}}(\gn{!E(x)},y)$ असल्यामुळे
\olref{repdiag2} वरून $y = \gn{!A}$ मिळते.
म्हणून $!B(y)$ वरून~$!B(\gn{!A})$ मिळते.

आता $!B(\gn{!A})$ मानू. \olref{repdiag1}
वरून पुढील मिळते:
\begin{align*}
& !D_{\fn{diag}}(\gn{!E(x)}, \gn{!A}) \land !B(\gn{!A}).
\intertext{त्यावरून पुढील निष्पन्न होते:}
& \lexists[y][(!D_{\fn{diag}}(\gn{!E(x)},y) \land !B(y))].
\end{align*}
पण हेच $!E(\gn{!E(x)})$, म्हणजेच~$!A$, आहे.
\end{proof}

\begin{digress}
या सिद्धतेची संगणनक्षमता सिद्धांतातील स्थिरबिंदू
उपपत्तीच्या सिद्धतेशी तुलना करा. येथे स्वतःच्याच
साहाय्याने \emph{विधान} परिभाषित करायचे आहे,
तर तिथे स्वतःच्याच साहाय्याने \emph{फलन}
परिभाषित करायचे होते. हा फरक सोडला, तर
मूळ कल्पना तीच आहे.
\end{digress}

\begin{prob}
  !!^a{formula}~$!A(x)$ ला \emph{सत्याची व्याख्या}
  असे म्हणतात, जर सर्व !!{sentence}s~$!B$
  साठी $\Th{Q}
  \Proves !B \liff !A(\gn{!B})$ असेल.
  स्थिरबिंदू उपपत्ती वापरून कोणतेही
  !!{formula} सत्याची व्याख्या असू शकत नाही
  हे दाखवा.
\end{prob}

\end{document}
